\documentclass[../../main.tex]{subfiles}
\begin{document}
% SOURCE p=246,247; printed=231-232; object=ch07-notes-title; region=bottom-top
\section{注记与参考文献}\label{section:ch07-notes-references}\label{section:ch07-notes}
% SOURCE p=246,247; printed=231-232; object=ch07-notes-riesz-history; region=bottom-upper
\paragraph*{\cref{section:ch07-positive-functionals}的注记。}
Riesz 表示定理的发展凝聚了许多人的工作。F.~Riesz\cite{folland-bib-116}首先证明了 $X=[a,b]$ 的情形；他的表述使用 Riemann--Stieltjes 积分，没有用到测度论。随后，Radon\cite{folland-bib-111}将它推广到 $\mathbb R^n$ 的紧子集，Banach 将它推广到紧度量空间（见 Saks\cite{folland-bib-128}），Kakutani\cite{folland-bib-080}将它推广到紧 Hausdorff 空间。在非紧情形，最早的一般结果由 Markov\cite{folland-bib-098}得到：他用有限可加集函数刻画了正规空间 $X$ 上 $BC(X)$ 的正线性泛函。Bourbaki 在约 1940 年似乎已经知道一个实质上等价于\cref{theorem:ch07-riesz-positive}的定理（见 Weil\cite[\S6]{folland-bib-158}），但其积分论著作直到 1952 年才出版；到那时，其他一些人也得到了类似结果。更详细的文献线索见 Dunford 与 Schwartz\cite[\S IV.16]{folland-bib-035}及 Hewitt 与 Ross\cite[\S11]{folland-bib-075}。关于向非局部紧空间推广 Riesz 表示定理，还可参见 König\cite{folland-bib-086}。

% SOURCE p=247; printed=232; object=ch07-notes-radon-terminology; region=upper-middle
“Radon 测度”一词源于 Radon 的奠基性论文\cite{folland-bib-111}；本书的用法很常见，却并非完全统一。有些作者将这样的测度称作“正则 Borel 测度”；另一些作者用“Radon 测度”指 $C_c(X)$ 上的正线性泛函；还有一些作者要求它在全部 Borel 集上内正则，而不是外正则。还须注意，一些较早的著作把 Borel $\sigma$ 代数定义为由紧集生成的 $\sigma$ 代数，它一般比本书的 $\mathscr B_X$ 小。
% SUPPLEMENT: id=v14-ch07-notes-terminology-details; parent=ch07-notes-radon-terminology; kind=omitted-proof+motivation+hypotheses; origin=independent-derivation
\begin{remark}
\textbf{不同约定的确切差别（重构新增）。}Hausdorff 空间的紧集闭，所以紧集生成的 $\sigma$ 代数包含于 $\mathscr B_X$，但未必等于它。例如在不可数离散空间中，紧集恰为有限集：无限集由各单点构成的开覆盖没有有限子覆盖，有限集则显然紧。因此紧集生成的 $\sigma$ 代数恰为可数或余可数集族；而每个子集都开，所以 Borel 代数是全部幂集。将该空间分成两个不可数子集，就得到 Borel 而不在前一代数中的集合。在第二可数 LCH 空间中，每个开集是可数个紧集的并，两种约定则相同：对开集内每点选闭包仍在该开集中的相对紧邻域，再取可数子覆盖即可。

这些术语差异影响定理的适用范围。正泛函与本书 Radon 测度之间的对应已经由 Riesz 定理建立，但把泛函称为测度仍是一项术语约定。要求全部 Borel 集内正则的版本不能与本书的定义直接互换；\cref{exercise:ch07-uncountable-lines}已给出本书 Radon 测度在某个无限测度集上不内正则的例子。查阅文献时须先核对定义及可测集合的范围。
\end{remark}
% END SUPPLEMENT: id=v14-ch07-notes-terminology-details

% SOURCE p=247; printed=232; object=ch07-notes-decomposable-extension; region=middle
若 $\mu$ 是 Radon 测度，以 $\overline\mu$ 表示\cref{section:ch07-positive-functionals}末尾讨论的完备、饱和延拓。一个重要事实是，$\overline\mu$ 总是\cref{exercise:ch03-decomposable-measure}意义下的可分解测度；见 Hewitt 与 Stromberg\cite[定理 19.30]{folland-bib-076}。因此可以使用该题中的 Radon--Nikodym 定理推广，以及\cref{exercise:ch06-linfty-semifinite-25}中的自然等距同构
\begin{align*}
(L^1(\overline\mu))^*\cong\mathcal L^\infty(\overline\mu).
\end{align*}
在这个联系中还应注意，由\cref{proposition:ch02-measurable-version-completion,proposition:ch02-finite-integral-support}，$p<\infty$ 时 $L^p(\overline\mu)$ 与 $L^p(\mu)$ 实质上相同。
% SUPPLEMENT: id=v14-ch07-notes-decomposability-proof; parent=ch07-notes-decomposable-extension; kind=omitted-proof+motivation+hypotheses; origin=independent-derivation
\begin{remark}
\textbf{可分解性的完整补证（重构新增）。}记延拓的定义域为 $\overline{\mathscr B}$。先取一个相对紧开覆盖，并将其良序排列为 $\{U_\alpha:\alpha<\kappa\}$。每个 $U_\alpha$ 测度有限。令
\begin{align*}
V_\alpha=\bigcup_{\beta<\alpha}U_\beta\quad(\alpha\leqslant\kappa),
\qquad D_\alpha=U_\alpha\setminus V_\alpha\quad(\alpha<\kappa).
\end{align*}
所有 $V_\alpha$ 开，所有 $D_\alpha$ 为有限测度 Borel 集；删去空集后，$D_\alpha$ 两两不交且并为 $X$。这已给出可分解定义的前两项。下面的关键是验证有限测度集不会把正质量“分散”到不可数多个零交块中。

固定有限测度 Borel 集 $E$，令 $a_\alpha=\mu(E\cap V_\alpha)$。若 $\lambda\leqslant\kappa$ 为极限序数，则
\begin{align*}
a_\lambda=\sup_{\alpha<\lambda}a_\alpha.
\end{align*}
确实，$E\cap V_\lambda$ 为有限测度 Borel 集，故内正则。它的每个紧子集 $K$ 被 $\{U_\beta:\beta<\lambda\}$ 覆盖，紧性给有限子覆盖；有限个下标在某个 $\alpha<\lambda$ 以前，所以 $K\subseteq V_\alpha$。对 $K$ 取上确界即得等式。又有 $a_0=0$ 及
\begin{align*}
a_{\alpha+1}-a_\alpha=\mu(E\cap D_\alpha).
\end{align*}
由超限归纳，$a_\gamma=\sum_{\alpha<\gamma}\mu(E\cap D_\alpha)$：后继步用刚才的差式；极限步用连续性，并注意任意有限个下标均在某个更早初段内，故非负族的有限部分和上确界也相同。在 $\gamma=\kappa$ 处，$V_\kappa=X$，得到
\begin{align*}
\mu(E)=\sum_{\alpha<\kappa}\mu(E\cap D_\alpha).
\end{align*}
右侧正项至多可数，因为有限和中大于 $1/n$ 的项只能有限。令 $J_E=\{\alpha:\mu(E\cap D_\alpha)>0\}$，则
\begin{align*}
N_E=E\setminus\bigcup_{\alpha\in J_E}D_\alpha
\end{align*}
是 Borel 零集：可数可加性和前式使被取走部分已有 $E$ 的全部质量。

若 $A\in\overline{\mathscr B}$ 且 $\overline\mu(A)<\infty$，\cref{section:ch07-positive-functionals}末尾的外覆盖论证说明 $A$ 属于原测度的完备化：存在 Borel 集 $E\subseteq A\subseteq B$，使 $\mu(B\setminus E)=0$、$\mu(B)<\infty$。因此 $\overline\mu(A)=\mu(E)$，且每个 $\alpha$ 都有 $\overline\mu(A\cap D_\alpha)=\mu(E\cap D_\alpha)$。上式于是给出可分解定义的第三项。

最后，设任意子集 $H\subseteq X$ 满足每个 $H\cap D_\alpha\in\overline{\mathscr B}$。对任意有限测度 Borel 集 $E$，将 $H\cap E$ 分成
\begin{align*}
H\cap E=
\bigcup_{\alpha\in J_E}(H\cap E\cap D_\alpha)\ \cup\ (H\cap N_E).
\end{align*}
前面是可数并；其中每项为有限测度的 $\overline{\mathscr B}$ 集，故在原完备化内。后一项为 Borel 零集的子集，也在完备化内。因此 $H\cap E$ 完备化可测。任意有限测度的完备化集合都有有限 Borel 上包集，所以同样结论对它也成立。于是 $H$ 关于完备化局部可测；前述饱和化刻画给 $H\in\overline{\mathscr B}$，验证第四项，证明完毕。$X=\varnothing$ 时使用空分解，四项直接成立。

动机是以良序开覆盖构造不交有限块，而用紧集内逼近保证极限序数处没有质量跳跃。LCH 给相对紧开覆盖，Radon 性给紧集有限和有限 Borel 集内正则，完备性处理零集的任意子集，饱和性完成逐块可测的拼接。良序选择沿用本书选择公理背景。这里只对\emph{有限}测度集证明第三项，未擅自把它加强到所有集合；因此并未推出 $\overline\mu$ 一定半有限。
\end{remark}
% END SUPPLEMENT: id=v14-ch07-notes-decomposability-proof
% SUPPLEMENT: id=v14-ch07-notes-lp-identification; parent=ch07-notes-decomposable-extension; kind=omitted-proof+motivation+hypotheses; origin=independent-derivation
\begin{remark}
\textbf{$L^p$ 的识别与对偶记号（重构新增）。}对 $0<p<\infty$，若 $f\in L^p(\overline\mu)$，则 $A_n=\{|f|>1/n\}$ 有有限测度，因为 $\overline\mu(A_n)\leqslant n^p\int|f|^p\,\mathrm d\overline\mu$。各 $A_n$ 有有限测度 Borel 上包集 $B_n$，令 $B=\bigcup_nB_n$。于是 $B$ 为 $\mu$ 的 $\sigma$ 有限 Borel 集，且 $f$ 在 $B^c$ 上为零。在 $B$ 上，饱和完备延拓与普通完备化相同；由\cref{proposition:ch02-measurable-version-completion}，$f|_B$ 有 Borel 可测且几乎处处相等的版本 $g_B$。在 $B^c$ 上补零得到 Borel 函数 $g$，与 $f$ 在 $\overline\mu$ 下几乎处处相等。简单函数逼近及单调收敛又表明，对任意 Borel 非负函数，其关于延拓与原测度的积分相同。因此 $\int|g|^p\,\mathrm d\mu=\int|f|^p\,\mathrm d\overline\mu$。反向嵌入显然也保持此积分；对两个 Borel 代表，差集的零测性在延拓前后不变，所以确为等距的一一对应（$p<1$ 时按相应准范数理解）。

这一步的动机是：有限 $p$ 的可积函数自动集中在可数个有限测度块上，因此新增的无限测度可测集不增加 $L^p$ 元素。$p=\infty$ 时没有这种有限水平集保证，不可照搬此证明。原文的 $\mathcal L^\infty$ 是\cref{exercise:ch06-linfty-semifinite-23}按\emph{局部零集}识别的空间；它不能在一般非半有限情形下不加说明地换成通常的 $L^\infty$。可分解性已验证\cref{exercise:ch06-linfty-semifinite-25}的全部假设，该题逐块表示及范数等式便给上面的对偶同构。
\end{remark}
% END SUPPLEMENT: id=v14-ch07-notes-lp-identification
\subfile{section-05-constructions.tex}
\end{document}
